% SPDX-FileCopyrightText: 2026 Christoph Maier
% SPDX-License-Identifier: Apache-2.0
%
% Proposed revision of the class-AB pad driver (IEEEtran version of proposed_improvements.md).
% Build: pdflatex proposed_improvements.tex (twice). Needs IEEEtran.cls (TeX Live: texlive-publishers);
% otherwise only amsmath, booktabs, tabularx, array, cite, textcomp, url.
\documentclass[journal]{IEEEtran}

\usepackage{amsmath}
\usepackage{booktabs,tabularx,array}
\usepackage{textcomp}
\usepackage[hyphens]{url}
\usepackage{cite}

\newcolumntype{L}{>{\raggedright\arraybackslash}X}
\newcommand{\cell}[1]{\texttt{#1}}
\newcommand{\um}{\,\textmu m}
\newcommand{\umsq}{\,\textmu m$^2$}
\newcommand{\uA}{\,\textmu A}
\newcommand{\uS}{\,\textmu S}
\newcommand{\mV}{\,mV}
\newcommand{\dB}{\,dB}
\newcommand{\MHz}{\,MHz}
\newcommand{\pF}{\,pF}
\newcommand{\kOhm}{\,k$\Omega$}
\newcommand{\Ohm}{\,$\Omega$}
\newcommand{\dg}{$^\circ$}
\newcommand{\RL}{1\,k$\Omega\parallel$100\,pF}
\newcommand{\vinp}{v_\mathrm{inp}}
\newcommand{\vinn}{v_\mathrm{inn}}
\newcommand{\vref}{v_\mathrm{ref}}
\newcommand{\vfb}{v_\mathrm{fb}}
\newcommand{\vout}{v_\mathrm{out}}

\begin{document}

\title{Proposed Revision of a Class-AB Analog Pad Driver in SG13CMOS5L: Matched-Pair DDA,
MOS Miller Capacitors, and a Resized Front End}

\author{C.~H.~Maier%
\thanks{Design note in the IEEE journal format; companion to
\cell{class\_ab\_pad\_driver.md}. Simulation only (ngspice-42, IHP-Open-PDK \cell{dev}, tt,
27\,\textdegree C unless stated), no layout, not reviewed. Drafted with Claude (configured
model \cell{claude-opus-5-5}; the serving model may differ), 2~October 2026. Running log:
\cell{log.md}; sources: \cell{literature\_notes.md}; every number:
\cell{sim/results\_improvements.txt}, produced by \cell{sim/run\_improvements.py}.}}

\markboth{Design note, sg13cmos5l\_cm\_ip\_\_single2diff2single, 2 October 2026}%
{Maier: Proposed Revision of a Class-AB Analog Pad Driver}

\maketitle

\begin{abstract}
The first version of a differential-to-single-ended class-AB pad driver in IHP's SG13CMOS5L
used the pad's ESD clamp frames as output devices. Its area was dominated by two MOM Miller
capacitors (1922 of 2387\umsq{} outside the frames). Its gain was set by a ratio of two
degenerated pairs, $(R_2+2/g_m)/(R_1+2/g_m)$, so it moved by $-2.6\ldots+5.0$\,\% over
corners. Its offset came from the folding sinks, and its loop gain was limited by short
internal cascodes. The poly resistors themselves occupy 28\umsq{}. The revision keeps the
frames and the floating class-AB control. It replaces the differential difference amplifier
(DDA) with four identical units balanced as $f(\vinp-\vref)+f(\vref-\vinn)=2f(\vfb-\vref)$,
which sets the gain $\tfrac12$ by the unit count and cancels the units' nonlinearity as far
as they match: a static ``externally linear, internally nonlinear'' circuit. The Miller
capacitors become thick-oxide PMOS accumulation capacitors, and the front end and the
internal cascodes are resized. Into 1\,k$\Omega\parallel$100\,pF the supply current drops
from 427 to 296\uA{}, the loop gain rises from 66.5 to 78.9\dB, and the gain over 11 corners
stays within 0.4992\ldots0.5000. The DC nonlinearity drops from 5.5 to 1.2\mV, THD at 10\,kHz
from $-51$ to $-58.5$\dB, the mismatch offset $\sigma$ from 14.4 to 5.3\mV, and the output
noise at 1\,kHz from 1.6 to 0.68\,\textmu V/$\sqrt{\mathrm{Hz}}$. Three resistor-free
degenerations from the literature were checked: a triode MOS, a well input, and an
undegenerated square-law pair. None fits this topology, because each unit's output must
depend on the gate difference alone. The output devices cannot be cascoded while they are
the ESD clamps. Moved into the slot and cascoded, they give a capacitor-free,
load-compensated driver with 74\dB{} of loop gain, for capacitive loads only.
\end{abstract}

\begin{IEEEkeywords}
Class-AB output stage, differential difference amplifier, externally linear circuits, ESD,
matched-nonlinearity cancellation, MOS capacitor, Miller compensation, pad driver, SG13CMOS5L.
\end{IEEEkeywords}

\section{Where the First Version's Numbers Come From}

\IEEEPARstart{T}{able}~\ref{tab:diag} traces each figure of merit of the first version
\cite{companion} to its cause. Only one of the causes is a linear passive: the MOM
capacitors. The resistors set the gain error, but through their combination with $g_m$, not
by being resistors.

\begin{table}[!t]
\caption{First Version: Figures of Merit and Their Causes}
\label{tab:diag}
\centering\footnotesize
\begin{tabularx}{\columnwidth}{@{}>{\raggedright\arraybackslash}p{0.2\columnwidth}>{\raggedright\arraybackslash}p{0.26\columnwidth}L@{}}
\toprule
quantity & value & traced to\\
\midrule
area outside frames & 2387\umsq{} (gate 437, MOM 1922, rhigh 28) & MOM Miller capacitors on M1--M4\\
supply current & 427\uA & output $I_Q$ 240, bias fixture $\approx$\,90, front end $\approx$\,100\\
gain over corners & 0.4870\ldots0.5249 & $(R_2+2/g_m)/(R_1+2/g_m)$: $R$ tracks $R$, $g_m$ does not\\
DC nonlinearity & 5.5\mV{} (12.7\mV{} at ss 27\,\textdegree C) & degenerated pairs, outside the loop\\
offset $\sigma$ (30 seeds) & 14.4\mV & sinks 11.5, tails 4.2, mirror 2.3, pairs 1.7\mV\textsuperscript{a}\\
loop gain, \RL & 66.5\dB & $A_1$ (short cascodes) $\times$ $A_2=g_{m,\mathrm{out}}(R_L\parallel r_o)\approx 6$\\
\bottomrule
\end{tabularx}\\[2pt]
\parbox{\columnwidth}{\footnotesize \textsuperscript{a}Budget of the MP-DDA in the first
version's sizing (fold, sinks and mirror identical), one device group at a time, 20 seeds.}
\end{table}

The offset budget follows Sarpeshkar's rule \cite[\S3.3.4]{sarpeshkar97}: the input offset
is the current mismatch divided by the transconductance,
\begin{equation}
V_{os} = \Delta I / G_m ,
\end{equation}
so a front end with a low $G_m/I$ is offset-prone, and the mirrors and sinks matter more than
the input devices.

\section{Matched-Pair DDA}

\subsection{Principle}
Four identical transconductance units each deliver $f(g_p-g_n)$ to the fold nodes:
unit~A with $(g_p,g_n)=(\vinp,\vref)$, unit~B with $(\vref,\vinn)$, and units C1 and C2,
both with $(\vref,\vfb)$, subtracting. The loop forces
\begin{equation}
f(\vinp-\vref)+f(\vref-\vinn)=2f(\vfb-\vref).
\label{eq:balance}
\end{equation}
With the input common mode at $\vref$, $\vinp-\vref=\vref-\vinn=u$, and (\ref{eq:balance})
reads $2f(u)=2f(\vout-\vref)$. Hence $\vout-\vref=(\vinp-\vinn)/2$ for any odd, monotonic $f$.
For an input common-mode error $c$, the input units carry
\begin{equation}
f(u+c)+f(u-c)=2f(u)+f''(u)\,c^2+O(c^4),
\end{equation}
so the error, $f''(u)\,c^2/2f'(u)$, is second order in $c$. The gain $\tfrac12$ is the ratio
$2{:}4$ of unit counts. The input units compress, and the feedback units apply the inverse of
the same law. This is the static form of ELIN: Minch's ``function inversion by negative
feedback around a high-gain amplifier'' \cite{minch97}, and the cancellation the
Säckinger--Guggenbühl DDA \cite{sackinger87} shows when both pairs carry the same signal.

\subsection{Requirement on a Unit}
One input of every unit sits at $\vref$. For the same differential voltage $x$, the input
common mode, and with it the source voltages, moves by $+x/2$ in A and C and by $-x/2$ in B.
So a unit qualifies only if its output depends on $g_p-g_n$ alone, not on where its sources
sit; its own linearity is secondary. The test bench sweeps
$f_A(x)=I(\vref+x,\vref)$ and $f_B(x)=I(\vref,\vref-x)$ and solves
$f_A(u)+f_B(u)=2f_A(v)$ (Table~\ref{tab:units}). For the resistor unit this gives
$-1.04/-1.02$\mV{} at $u=\mp0.5$\,V. The full driver's largest deviation from its
small-signal line is 1.04\mV.

\begin{table*}[!t]
\caption{DDA Units: Common-Mode Sensitivity and Static MP-DDA Error}
\label{tab:units}
\centering\footnotesize
\begin{tabular}{@{}lcccc@{}}
\toprule
unit & $G_{m0}$ & $G_m$ vs.\ input CM & own compression at 0.5\,V & DDA error at $u=-0.5/+0.5$\,V\\
\midrule
\cell{unit\_r}: split tails, rhigh between the sources & 15.6\uS & 0.6\,\%/V & 0.5\,\% & $-1.0/-1.0$\mV\\
\cell{unit\_t}: same, triode hv-PMOS instead of rhigh & 15.5\uS & 77\,\%/V & 20\,\% & $-191/-72$\mV\\
\cell{unit\_w}: well (bulk) input, gates at $v_{ss}$ \cite{sarpeshkar97} & 11.0\uS & $-27$\,\%/V & 6.8\,\% & $+29/+38$\mV\\
\cell{unit\_q}: undegenerated square-law pair, 2/6\um & 7.8\uS & $-72$\,\%/V & 22\,\% & $+75/+200$\mV\\
\bottomrule
\end{tabular}
\end{table*}

The split-tail resistor unit qualifies because each input device is a constant-current
source follower. A follower's transfer does not depend on the device law
\cite[\S7.2.1]{minch97}, so the element between the sources sees exactly $g_p-g_n$. The other
three fail for one common reason. A triode MOS's resistance depends on its $V_{SG}$, which is
the absolute source voltage. The well input's $\kappa$ depends on the well-to-gate voltage
\cite[\S3.3.1]{sarpeshkar97}. A single-tail pair's tail and channel-length modulation see the
source move. The design rule is therefore ``no element whose value depends on the source
voltage'', not ``no resistor'', and the poly resistor stays. A unit cannot be made symmetric
in common mode here, because the feedback side has only $\vfb$ and $\vref$; that would need
an inverted replica of the output.

\subsection{Result With Nothing Else Changed}
\cell{d2s\_mp} keeps the first version's fold, mirror, class-AB control, frames, MOM
capacitors and bias, and the same 80\uA{} of tail current as $8\times10$\uA{}
(Table~\ref{tab:mp}). The gain $\sigma$ is now the rhigh unit mismatch: 1.0\,\% for
$0.5\times36.4$\um{} units and 0.84\,\% for the $0.5\times50.6$\um{} units of
\cell{d2s\_mpdda}. Four times the resistor area would halve it, at about 75\umsq{} more per
unit. The offset is untouched, because it lives in the fold (Section~\ref{sec:sizing}).

\begin{table}[!t]
\caption{MP-DDA in the First Version's Sizing}
\label{tab:mp}
\centering\footnotesize
\begin{tabularx}{\columnwidth}{@{}Lcc@{}}
\toprule
 & first version & MP-DDA\\
\midrule
gain, tt & 0.5005 & 0.4998\\
gain, 11 corners & 0.4870\ldots0.5249 & 0.4983\ldots0.4998\\
nonlinearity $\pm1$\,V, tt / worst & 5.52 / 12.69\mV & 1.04 / 1.91\mV\\
THD 10/100\,kHz, \RL & $-51.1/-50.2$\dB & $-60.2/-55.5$\dB\\
THD 10/100\,kHz, 100\pF & $-51.0/-51.0$\dB & $-60.0/-59.8$\dB\\
$T_0$ (\RL) & 66.5\dB & 66.5\dB\\
$f_c$ / PM (\RL) & 1.62\MHz{} / 74.8\dg & 1.62\MHz{} / 74.7\dg\\
gain $\sigma$ / offset $\sigma$ (30 seeds) & 1.60\,\% / 14.4\mV & 1.33\,\% / 14.5\mV\\
\bottomrule
\end{tabularx}
\end{table}

\section{Miller Compensation Without BEOL Capacitors}

\subsection{Device}
Each Miller capacitor is an \cell{sg13\_hv\_pmos} used as an accumulation capacitor, with
the gate above the n-well. On the A side the gate is node $a$ ($\approx$\,2.6\,V) and
source, drain and well are $\vout$. On the B side the gate is $\vout$ and source, drain and
well are node $b$ ($\approx$\,0.7\,V). Over the linear output range $V_{GW}$ spans
0.15\ldots1.15\,V (A) and 0.25\ldots1.25\,V (B). $C_{ox}\approx5$\,fF/\textmu m$^2$
($t_{ox}=6.9$\,nm), and in that range the device delivers 0.38\ldots0.76 of $C_{ox}A$. A
$14\times14$\um{} device has 0.370\pF{} at 0.15\,V, 0.648\pF{} at 1.0\,V and 0.745\pF{} at
1.25\,V; a $31\times31$\um{} MOM capacitor has 1.002\pF.

\subsection{Equivalence}
$18\times18$\um{} per side reproduces the MOM pair's loop within a degree
(Table~\ref{tab:cl}). That is 648\umsq{} of gate oxide instead of 1922\umsq{} of M1--M4
fingers. THD with capacitive loads is unchanged ($-60.0/-60.0$\dB{} at 100\pF). With \RL{},
100\,kHz THD loses 2.3\dB{} ($-53.2$ vs.\ $-55.5$\dB): the C--V curvature modulates $f_c$
where the loop gain is lowest.

\begin{table}[!t]
\caption{Crossover Frequency / Phase Margin vs.\ Load Capacitance (No Resistive Load)}
\label{tab:cl}
\centering\footnotesize
\begin{tabular}{@{}ccc@{}}
\toprule
$C_L$ & MOM $31\times31$\um & MOS $18\times18$\um\\
\midrule
5\pF & 2.18\MHz{} / 86.1\dg & 2.14\MHz{} / 86.4\dg\\
20\pF & 2.17\MHz{} / 82.6\dg & 2.14\MHz{} / 82.9\dg\\
100\pF & 2.03\MHz{} / 66.2\dg & 2.00\MHz{} / 67.0\dg\\
300\pF & 1.60\MHz{} / 45.5\dg & 1.59\MHz{} / 46.3\dg\\
1\,nF & 1.00\MHz{} / 26.4\dg & 1.00\MHz{} / 26.9\dg\\
\bottomrule
\end{tabular}
\end{table}

\subsection{Why the Capacitor Cannot Simply Be Made Small}
With Miller compensation the non-dominant pole is
\begin{equation}
\omega_{p2}\approx\frac{g_{m,\mathrm{out}}\,C_M}{C_L\,(C_g+C_M)} .
\end{equation}
The frame gates carry $C_g\approx0.86$\pF{} (P) and 0.69\pF{} (N), so pole splitting needs
$C_M$ of the order of $C_g$ whatever $G_{m1}$ is; lowering $G_{m1}$ only lowers $f_c$. Two
ways around this floor were considered but not simulated:
\begin{itemize}
\item Source-follower buffers in front of the frame gates. They shrink the high-impedance
node, so $C_M$ could drop to perhaps 0.2\pF{} (estimate), but the class-AB loop would have
to include the followers.
\item Indirect (current-buffer) compensation \cite{ahuja83} into the fold nodes. These sit
at $\approx$\,0.37\,V, which biases a MOS capacitor poorly. The $a$-side current would also
have the wrong sign through the mirror unless injected at the mirror cascode's source.
\end{itemize}
Load compensation needs no capacitor at all, but only once the output can be cascoded
(Section~\ref{sec:casc}).

\subsection{ESD Check Item}
The A-side capacitor puts a p$^+$/n-well junction, the n-well/p-substrate diode and 7\,nm of
gate oxide between the pad and node $a$. The B-side capacitor puts gate oxide between the
pad and node $b$. The frames' own gate--drain overlaps already carry the same oxide stress,
but the capacitors belong in the ESD review.

\section{Currents, Offset and Noise: \texttt{d2s\_mpdda}}
\label{sec:sizing}

\begin{table*}[!t]
\caption{ELIN Approaches With Non-Exponential Internal Laws: Feasibility for This Driver}
\label{tab:elin}
\centering\footnotesize
\begin{tabularx}{\textwidth}{@{}>{\raggedright\arraybackslash}p{0.2\textwidth}>{\raggedright\arraybackslash}p{0.17\textwidth}>{\raggedright\arraybackslash}p{0.2\textwidth}L@{}}
\toprule
approach & internal law & sources & verdict\\
\midrule
matched-nonlinearity feedback, $f^{-1}\circ f$ & any odd monotonic $f$ (here a degenerated square law) & \cite{minch97,sackinger87} & \textbf{used}; the unit must depend on $g_p-g_n$ only (Section~II)\\
MOS translinear loop & square law & \cite{wiegerink,seevinck87} & already present: the floating class-AB control \cite{hogervorst94,monticelli86} is one\\
constant-sum square-law push-pull & square law & \cite{sarpeshkar97} & the frames run at $g_m/I\approx12.5$\,V$^{-1}$ (moderate inversion); an exact square law would need mA of $I_Q$; not pursued\\
$\sqrt{\ }$-domain / dynamic translinear & square law or exponential & \cite{mulder97tcas,mulder97jssc,efthivoulidis99,hiseni09} & dynamic (filters); a static buffer has nothing to compand\\
triode MOSFET-C & triode, even orders cancel & --- & needs balanced signals; the feedback side here is single-ended\\
capacitive dividers: FG input, MITE, AFGA, MIFG & capacitive + exponential or law-independent follower & \cite{minch97,minch01,hasler97,sanchez_fg,montalvo93} & needs linear capacitors, charge control (tunneling, injection or UV) and AC coupling; none is available or acceptable here\\
well input, bump linearization & $\kappa$-scaled law; bump in subthreshold & \cite{sarpeshkar97} & $G_m$ depends on input CM: $+29/+38$\mV{} error (\cell{unit\_w}); bump range is subthreshold-sized\\
undegenerated strong-inversion pair & square law & --- & linear enough, but the sources move with the signal: $+75/+200$\mV{} error (\cell{unit\_q})\\
\bottomrule
\end{tabularx}
\end{table*}

\begin{table*}[!t]
\caption{Variants Compared (\RL{} Unless Stated)}
\label{tab:opt}
\centering\footnotesize
\begin{tabular}{@{}lcccccc@{}}
\toprule
variant & area & $I_{DD}$ & offset $\sigma$ & $T_0$: \RL{} / 50\Ohm{} / 100\pF & gain over corners & THD 10/100\,kHz\\
\midrule
first version (\cell{d2s\_miller}) & 2387\umsq & 427\uA & 14.4\mV & 66.5/40.7/96.5\dB & $-2.6\ldots+5.0$\,\% & $-51.1/-50.2$\dB\\
area-lean: \cell{d2s\_mp} + MOS $18\times18$ & 1238\umsq & 427\uA & 14.5\mV & 66.5/40.7/96.5\dB & $-0.3\ldots0$\,\% & $-60.1/-53.2$\dB\\
\cell{d2s\_mpdda}, $L_\mathrm{cas}=1$\um & 1694\umsq & 296\uA & 5.3\mV\textsuperscript{a} & 69.7/43.9/99.6\dB & not simulated & $-58.7/-51.4$\dB\textsuperscript{a}\\
\cell{d2s\_mpdda}, $L_\mathrm{cas}=3$\um{} (default) & 2014\umsq & 296\uA & 5.3\mV & 78.9/53.1/104.3\dB & $-0.16\ldots0$\,\% & $-58.5/-51.3$\dB\\
\bottomrule
\end{tabular}\\[2pt]
\parbox{\textwidth}{\footnotesize \textsuperscript{a}From the intermediate run in
\cell{log.md} (\cell{d2s\_mp2}: the circuit of \cell{d2s\_mpdda} with $L_\mathrm{cas}=1$\um{}
except $n_g=1$ in the cascodes; its loop numbers agree with the $L_\mathrm{cas}=1$ row of
\cell{results\_improvements.txt}), not from that file itself.}
\end{table*}

Headroom fixes the sizing. Nodes $a$ ($\approx$\,2.6\,V) and $b$ ($\approx$\,0.7\,V) are the
frame gates, which leaves about 0.35\,V per device for the mirror and its cascode and for the
sinks and theirs. The overdrive of the matching-critical devices cannot be raised, so their
area is raised instead (Table~\ref{tab:sizing}). All conducting devices are at least 80\mV{}
clear of saturation at $v_d=0$. Table~\ref{tab:perf} compares the result with the first
version.

\begin{table}[!t]
\caption{Sizing Changes}
\label{tab:sizing}
\centering\footnotesize
\begin{tabularx}{\columnwidth}{@{}lLL@{}}
\toprule
block & first version & \cell{d2s\_mpdda}\\
\midrule
unit tails & 2 pairs $\times$ 2 $\times$ 20\uA{} = 80\uA & 4 units $\times$ 2 $\times$ 5\uA{} = 40\uA{} (5/6\um)\\
degeneration & rhigh 116\kOhm{} / 51\kOhm & rhigh 150\kOhm{} per unit ($I_tR=0.75$\,V)\\
folding sinks & $2\times50$\uA, 25/2 & $2\times30$\uA, 36/6\\
mirror, cascodes & 10/2, $L=1$\um{} cascodes & 24/4, $L=3$\um{} cascodes ($W=10L$)\\
bias references & 90\uA & 26\uA\\
output $I_Q$ (P/N) & 240\uA & 212\uA{} (class-AB devices 5.2/4.8\uA)\\
$I_{DD}$ & 427\uA & 296\uA\\
\bottomrule
\end{tabularx}
\end{table}

\begin{table}[!t]
\caption{Performance: First Version vs.\ \texttt{d2s\_mpdda}}
\label{tab:perf}
\centering\footnotesize
\begin{tabularx}{\columnwidth}{@{}Lcc@{}}
\toprule
 & first version & \cell{d2s\_mpdda}\\
\midrule
$T_0$: \RL{} / 50\Ohm{} / 100\pF & 66.5/40.7/96.5\dB & 78.9/53.1/104.3\dB\\
$f_c$/PM, 100\pF & 2.02\MHz/66.4\dg & 1.77\MHz/62.0\dg\\
$f_c$/PM, 300\pF & 1.59\MHz/45.6\dg & 1.39\MHz/42.0\dg\\
$f_c$/PM, 1\,nF & 1.00\MHz/26.5\dg & 0.87\MHz/24.1\dg\\
gain, 11 corners & 0.4870\ldots0.5249 & 0.4992\ldots0.5000\\
nonlinearity, tt / worst (ss 27\,\textdegree C) & 5.52/12.69\mV & 1.24/3.50\mV\\
THD 10/100\,kHz, \RL & $-51.1/-50.2$\dB & $-58.5/-51.3$\dB\\
THD 10/100\,kHz, 50\Ohm & $-46.2/-29.1$\dB & $-45.4/-26.9$\dB\\
$t_{1\%}$, \RL & 307\,ns & 327\,ns\\
offset $\sigma$ / gain $\sigma$ (30 seeds) & 14.4\mV{} / 1.60\,\% & 5.3\mV{} / 1.11\,\%\\
noise, 1\,kHz (nV/$\sqrt{\mathrm{Hz}}$) & 1599 & 676\\
noise, 100\,kHz (nV/$\sqrt{\mathrm{Hz}}$) & 200 & 152\\
noise, 10\,Hz--10\,MHz (\textmu V rms) & 278 & 231\\
\bottomrule
\end{tabularx}
\end{table}

The offset budget of \cell{d2s\_mpdda} ($\sigma$, one group at a time, 20 seeds) is:
sinks 4.3\mV, unit pairs 1.7\mV, tails 1.4\mV, mirror 1.4\mV, the rest below 0.1\mV. The
sinks still lead.

The halved front-end $G_m$ lowers $f_c$ from 1.62 to 1.37\MHz{} (\RL). That costs the
100\,kHz THD gain of Section~II and makes 50\Ohm{} slightly worse. $14\times14$\um{}
capacitors raise $f_c$ to 2.12\MHz{} at 100\pF{} but leave 54\dg{} of phase margin; the
choice belongs to the specification. 50\Ohm{} remains what the first note called it: a DC
and low-frequency case.

\section{Cascoding}
\label{sec:casc}

\subsection{Output Devices Are the Pad's ESD Clamps}
A series cascode would put its own drain, not the clamp's, on the pad. A stacked clamp frame
would be a new ESD structure needing its own qualification. The output therefore stays
uncascoded, and the frames' $r_o$ ($V_A\approx26$\,V for P; $L=1.0$\um{} for N) bounds the
second stage. What can be cascoded are the internal high-impedance nodes $a$ and $b$
(Table~\ref{tab:lcas}). Regulated (gain-boosted) cascodes do not fit: the $\approx$\,0.35\,V
per device next to $a$ and $b$ leaves no room for an auxiliary amplifier's $V_{GS}$.

\begin{table}[!t]
\caption{Loop Gain vs.\ Fold and Mirror Cascode Length ($W=10L$)}
\label{tab:lcas}
\centering\footnotesize\setlength{\tabcolsep}{4pt}
\begin{tabular}{@{}cccccc@{}}
\toprule
$L$ & $T_0$, \RL & $T_0$, 100\pF & $T_0$, 50\Ohm & PM, 100\pF & area\\
\midrule
1\um & 69.7\dB & 99.6\dB & 43.9\dB & 63.4\dg & 1081\umsq\\
2\um & 75.4\dB & 103.1\dB & 49.6\dB & 62.8\dg & 1201\umsq\\
3\um & 78.9\dB & 104.3\dB & 53.1\dB & 62.0\dg & 1401\umsq\\
\bottomrule
\end{tabular}
\end{table}

\subsection{Output Devices in the Slot}
With the pad reached through \cell{IOPadAnalog}, the output devices are ordinary devices and
can be cascoded. \cell{d2s\_lc2} is load-compensated and has no compensation capacitor. It
uses the MP-DDA front end and half-frame output devices, with the cascode gates at $\vref$
(output range $\approx$\,1.0\ldots2.4\,V); see Table~\ref{tab:slot}. Cascoding adds 37\dB.
It helps neither resistive loads, which collapse any single-stage loop, nor the minimum load
capacitance ($\approx$\,30--40\pF).

\begin{table}[!t]
\caption{In-Slot, Load-Compensated Driver Without Compensation Capacitor}
\label{tab:slot}
\centering\footnotesize
\begin{tabular}{@{}lcc@{}}
\toprule
 & uncascoded & cascoded\\
\midrule
$T_0$ & 37.2\dB & 74.0\dB\\
PM, 5\pF & 16\dg & unstable\\
PM, 20\pF & 47\dg & 39\dg\\
PM, 100\pF & 78\dg & 77\dg\\
PM, 1\,nF & 89\dg & 89\dg\\
\RL & no loop gain & no loop gain\\
$I_{DD}$ & 174\uA & 169\uA\\
\bottomrule
\end{tabular}
\end{table}

In \cell{IOPadAnalog}, \cell{padres} sits behind the 586.9\Ohm{} of
\cell{SecondaryProtection}. Sensing the loop at \cell{padres} (the near side) turns that
resistor into an isolation resistor. The phase margin rises from 39\dg{} to 62\dg{} at
20\pF{} and reaches 104\dg{} at 100\pF; the pad then sees a $587\,\Omega\cdot C_L$ low-pass
outside the loop, 2.7\MHz{} at 100\pF. Sensing at \cell{pad} (the far side) keeps the
margins of the bare stage. 50\Ohm{} cannot be driven through 587\Ohm.

Two points remain open. First, the cascoded version has 7.8\mV{} of systematic offset,
because the diode replicas sit at a different $V_{DS}$ than the cascoded output devices;
cascoding the replicas should remove it. Second, an in-slot output stage is not an area
saving. This is an estimate, not simulated: for the same output pole at the same current,
in-slot devices for the two-stage topology would need about the frames' $W$, because the
frames' $g_m/I\approx12.5$\,V$^{-1}$ at 240\uA{} is what sets the non-dominant pole.

\section{Externally Linear Circuits With a Non-Exponential Internal Law}
Edwards \cite[p.~72]{edwards_thesis} (see also \cite{edwards00,edwards98}) names the square law of strong inversion
\cite{wiegerink} and the exponential as the two nonlinearities that ELIN circuits can
exploit. Table~\ref{tab:elin} lists what else the sources offer and how each fares here.
Two expressions recur in it. Sarpeshkar's square-law identity \cite[p.~25]{sarpeshkar97} is
\begin{equation}
(x-a)^2-(x-b)^2=(b-a)(2x-a-b).
\end{equation}
For a push-pull stage with constant $V_{GS,P}+V_{GS,N}=2V_Q$ and drive $\pm\delta$, it gives
\begin{equation}
I_P-I_N=2\beta\,(V_Q-V_t)\,\delta,\qquad |I_\mathrm{out}|<4I_Q ,
\end{equation}
exactly, in strong inversion.



ELIN without the exponential is therefore feasible here in one form: function inversion by
matched units. Its precondition is not a particular law but a unit whose output is
independent of the source voltage, which the source-follower-plus-resistor unit provides.
The square-law translinear classics need either strong inversion at currents this driver
does not spend, or balanced signals that a single-ended pad does not have.

\section{Options}
Table~\ref{tab:opt} compares the variants. Area is drawn device area (W$\cdot$L, resistor
w$\cdot$l, capacitor w$\cdot$l) without the frames and the bias fixture, as in the first
note. The area-lean variant's corner gain is that of \cell{d2s\_mp}, since the capacitors do
not enter the DC gain.



\section{Not Done}
\begin{itemize}
\item xschem schematics and the CACE run for \cell{d2s\_mpdda}. The existing yaml should take
it as a new DUT, with \cell{d2s\_bias\_lp} as fixture.
\item Layout. The four units must match, so common-centroid placement and Edwards's ``same
surround'' rule \cite[\S3.10]{edwards_thesis} apply; the PMOS input devices sit in separate
wells (bulk at their own source).
\item PSRR, CMRR, input CM range, supply corners (3.0/3.6\,V) and process Monte Carlo
(\cell{tt\_stat}). Only mismatch Monte Carlo was run.
\item ESD review of the pad-connected MOS capacitors (Section~III).
\item \cell{d2s\_lc2}: cascoded replicas, a real minimum-$C_L$ specification, and whether a
slot-side driver is wanted at all.
\item Inherited from the first version: the PMOS of the left floating replica (\cell{XFPL})
carries 1\,nA, because with its bulk at $v_{dd}$, $|V_{GS}|=0.78$\,V is below its
body-shifted $|V_{th}|$ of 1.01\,V. Only \cell{XFNL} replicates the class-AB control. The
circuit works, but the replica is not a true replica.
\end{itemize}

\appendix[Files]
All files are in \cell{sim/}:
\begin{itemize}
\item \cell{d2s\_mpdda.spice}: proposed driver, subcircuit \cell{d2s\_mpdda} (pins as
\cell{d2s\_miller}; parameters \cell{wc}, \cell{lc}, \cell{lcas}, \cell{rl}).
\item \cell{units.spice}: DDA units \cell{unit\_r}, \cell{unit\_r2} (used by
\cell{d2s\_mpdda}), \cell{unit\_t}, \cell{unit\_w}, \cell{unit\_q}.
\item \cell{d2s\_bias\_lp.spice}: bias fixture for \cell{d2s\_mpdda} (pins as
\cell{d2s\_bias}).
\item \cell{d2s\_mp.spice}, \cell{ccomp.spice}: MP-DDA in the first version's sizing; MOM or
MOS capacitors through a test-bench wrapper.
\item \cell{d2s\_lc2.spice}: in-slot, load-compensated, cascoded (\cell{d2s\_lc2}) and
uncascoded (\cell{d2s\_lc2\_nc}).
\item \cell{run\_improvements.py}: reproduces every number here; imports
\cell{../../run\_d2s.py} unchanged.
\item \cell{results\_improvements.txt}: its output (3\,min 49\,s on two cores).
\end{itemize}

\begin{thebibliography}{99}
\bibitem{companion}
C.~H.~Maier, ``Two proof-of-concept class-AB analog pad drivers built on ESD-clamp output
devices in IHP SG13CMOS5L,'' design note \cell{class\_ab\_pad\_driver.md}, Sep.~2026.

\bibitem{sarpeshkar97}
R.~Sarpeshkar, ``Efficient precise computation with noisy components: Extrapolating from an
electronic cochlea to the brain,'' Ph.D. dissertation, California Inst. Technol., Pasadena,
CA, USA, 1997, ch.~3.

\bibitem{minch97}
B.~A.~Minch, ``Analysis, synthesis, and implementation of networks of multiple-input
translinear elements,'' Ph.D. dissertation, California Inst. Technol., Pasadena, CA, USA,
1997.

\bibitem{minch01}
B.~A.~Minch, ``Multiple-input translinear element log-domain filters,'' \emph{IEEE Trans.
Circuits Syst. II}, vol.~48, no.~1, pp.~29--36, Jan.~2001.

\bibitem{hasler97}
P.~Hasler, Ph.D. dissertation, California Inst. Technol., Pasadena, CA, USA, ch.~4--5
(local copy: \cell{FloatingGate/Hasler/thesis\_chapter*.ps}).

\bibitem{sanchez_fg}
E.~S\'anchez-Sinencio, ``Floating gate techniques and applications,'' lecture notes, Analog
and Mixed-Signal Center, Texas A\&M Univ., College Station, TX, USA.

\bibitem{montalvo93}
A.~J.~Montalvo and J.~J.~Paulos, ``Improved floating-gate devices using standard CMOS
technology,'' \emph{IEEE Electron Device Lett.}, vol.~14, no.~8, pp.~372--374, Aug.~1993.

\bibitem{edwards_thesis}
R.~T.~Edwards, Ph.D. dissertation, Dept. Elect. Comput. Eng., Johns Hopkins Univ.,
Baltimore, MD, USA, \S\S3.5--3.10 (local copy: \cell{TimEdwards/thesis.pdf}).

\bibitem{edwards00}
R.~T.~Edwards and G.~Cauwenberghs, ``Synthesis of log-domain filters from first-order
building blocks,'' \emph{Analog Integr. Circuits Signal Process.}, vol.~22, pp.~177--186,
2000.

\bibitem{edwards98}
R.~T.~Edwards and G.~Cauwenberghs, ``A second-order log-domain bandpass filter for audio
frequency applications,'' in \emph{Proc. IEEE ISCAS}, Monterey, CA, USA, 1998, vol.~3.

\bibitem{efthivoulidis99}
G.~Efthivoulidis and Y.~Tsividis, ``Signal analysis of externally linear filters,'' in
\emph{Proc. IEEE ISCAS}, 1999, vol.~6, pp.~65--68.

\bibitem{mulder97tcas}
J.~Mulder, A.~C.~van~der~Woerd, W.~A.~Serdijn, and A.~H.~M.~van~Roermund, ``General
current-mode analysis method for translinear filters,'' \emph{IEEE Trans. Circuits Syst. I},
vol.~44, no.~3, pp.~193--197, Mar.~1997.

\bibitem{mulder97jssc}
J.~Mulder, A.~C.~van~der~Woerd, W.~A.~Serdijn, and A.~H.~M.~van~Roermund, ``An RMS-DC
converter based on the dynamic translinear principle,'' \emph{IEEE J. Solid-State Circuits},
vol.~32, no.~7, pp.~1146--1150, Jul.~1997.

\bibitem{hiseni09}
S.~Hiseni, C.~Sawigun, and W.~A.~Serdijn, ``Dynamic translinear nonlinear energy operator,''
in \emph{Proc. ECCTD}, 2009, pp.~153--156.

\bibitem{wiegerink}
R.~J.~Wiegerink, \emph{Analysis and Synthesis of MOS Translinear Circuits}. The Hague, The
Netherlands (as cited in \cite{edwards_thesis}).

\bibitem{sackinger87}
E.~S\"ackinger and W.~Guggenb\"uhl, ``A versatile building block: The CMOS differential
difference amplifier,'' \emph{IEEE J. Solid-State Circuits}, vol.~SC-22, no.~2,
pp.~287--294, Apr.~1987.

\bibitem{seevinck87}
E.~Seevinck and R.~F.~Wassenaar, ``A versatile CMOS linear transconductor/square-law function
circuit,'' \emph{IEEE J. Solid-State Circuits}, vol.~SC-22, no.~3, pp.~366--377, Jun.~1987.

\bibitem{hogervorst94}
R.~Hogervorst, J.~P.~Tero, R.~G.~H.~Eschauzier, and J.~H.~Huijsing, ``A compact
power-efficient 3~V CMOS rail-to-rail input/output operational amplifier for VLSI cell
libraries,'' \emph{IEEE J. Solid-State Circuits}, vol.~29, no.~12, pp.~1505--1513,
Dec.~1994.

\bibitem{monticelli86}
J.~M.~Monticelli, ``A quad CMOS single-supply op amp with rail-to-rail output swing,''
\emph{IEEE J. Solid-State Circuits}, vol.~SC-21, no.~6, pp.~1026--1034, Dec.~1986.

\bibitem{ahuja83}
B.~K.~Ahuja, ``An improved frequency compensation technique for CMOS operational
amplifiers,'' \emph{IEEE J. Solid-State Circuits}, vol.~SC-18, no.~6, pp.~629--633,
Dec.~1983 (from memory, not re-checked).
\end{thebibliography}

\end{document}
